\documentclass{csdept}

% ============================================================
% Document metadata
% ============================================================

\university{University of El Oued}
\faculty{Faculty of Exact Sciences}
\department{Department of Computer Science}

\course{Algorithms and Data Structures}
\teacher{Dr. Firstname Lastname}
\academicyear{2026--2027}

\documenttype{Reference Document}

\title{Computer Science Department Document System}

% ============================================================
% Document
% ============================================================

\begin{document}

\maketitle

\tableofcontents

\newpage

% ============================================================
\section{Introduction}
% ============================================================

This document demonstrates the main features provided by the
Computer Science Department document system.

The system provides support for:

\begin{itemize}
    \item department and course metadata;
    \item mathematical notation;
    \item theorem-like environments;
    \item exercises and teaching boxes;
    \item pseudocode and algorithms;
    \item source-code listings;
    \item standard LaTeX lists and structures;
    \item hyperlinks and references.
\end{itemize}

The professor only needs to load the \texttt{csdept} document class:

\begin{verbatim}
\documentclass{csdept}
\end{verbatim}

The document metadata can then be specified using commands such as
\texttt{\textbackslash course}, \texttt{\textbackslash teacher},
and \texttt{\textbackslash documenttype}.

% ============================================================
\section{Mathematics}
% ============================================================

The system provides mathematical typesetting through
\texttt{mathtools}.

Inline mathematics can be written normally, for example
$a^2+b^2=c^2$.

Displayed mathematics can be written as:

\[
    f(x) = x^2 + 2x + 1.
\]

More complicated expressions are also supported:

\[
    \sum_{i=1}^{n} i
    =
    \frac{n(n+1)}{2}.
\]

An equation can be numbered using the standard \texttt{equation}
environment:

\begin{equation}
    T(n) = 2T\left(\frac{n}{2}\right) + O(n).
\end{equation}

Matrices are supported as well:

\[
    A =
    \begin{pmatrix}
        1 & 2 \\
        3 & 4
    \end{pmatrix}.
\]

% ============================================================
\section{Theorem Environments}
% ============================================================

The theorem package provides several standard environments.

\subsection{Definition}

\begin{definition}
A graph is an ordered pair $G=(V,E)$ where $V$ is a set of
vertices and $E$ is a set of edges connecting pairs of vertices.
\end{definition}

\subsection{Theorem}

\begin{theorem}
Every finite tree with $n$ vertices has exactly $n-1$ edges.
\end{theorem}

\begin{proof}
We prove the result by induction on $n$.

For $n=1$, the tree contains one vertex and no edges, so the
statement holds.

Assume that every tree with $n$ vertices has $n-1$ edges.
A tree with $n+1$ vertices contains a leaf. Removing this leaf
and its incident edge produces a tree with $n$ vertices.
By the induction hypothesis, the remaining tree has $n-1$ edges.
Adding the removed edge gives $n$ edges.

Therefore, a tree with $n+1$ vertices has $n$ edges.
\end{proof}

\subsection{Lemma}

\begin{lemma}
If a graph contains a path from $u$ to $v$, then $u$ and $v$
belong to the same connected component.
\end{lemma}

\subsection{Proposition}

\begin{proposition}
Every connected graph contains a spanning tree.
\end{proposition}

\subsection{Corollary}

\begin{corollary}
Every connected graph with $n$ vertices contains a spanning tree
with $n-1$ edges.
\end{corollary}

\subsection{Example}

\begin{example}
Consider the path

\[
    v_1 \rightarrow v_2 \rightarrow v_3 \rightarrow v_4.
\]

This graph is connected and contains a spanning tree consisting
of all four vertices and the three edges shown above.
\end{example}

\subsection{Remark}

\begin{remark}
The theorem-like environments are provided by \texttt{amsthm}.
The standard \texttt{proof} environment is therefore available
without defining a custom environment.
\end{remark}

% ============================================================
% Figures
% ============================================================

\section{Figures}

Figures can be included using the \texttt{graphicx} package.
Use the \texttt{figure} environment to add a caption and label
to an image, allowing it to be referenced elsewhere in the document.

\subsection{Including an Image}

\begin{figure}[htbp]
    \centering
    \includegraphics[width=0.6\textwidth]{reference/figures/dep_logo.png}
    \caption{An example figure.}
    \label{fig:example}
\end{figure}

Figure~\ref{fig:example} illustrates how to include an image
with a caption and a reference label.

\subsection{Controlling Image Size}

The width of an image can be specified relative to the text width.

\begin{figure}[htbp]
    \centering
    \includegraphics[width=0.35\textwidth]{reference/figures/dep_logo.png}
    \caption{An image displayed at a smaller size.}
    \label{fig:small}
\end{figure}

% ============================================================


% ============================================================
% Tables
% ============================================================

\section{Tables}

Tables can be created using the \texttt{tabular} environment.
The \texttt{booktabs} package provides professional horizontal
rules, while \texttt{tabularx} allows tables to fit the text width.

\subsection{Basic Table}

\begin{table}[htbp]
    \centering
    \caption{Algorithm complexity comparison.}
    \label{tab:complexity}

    \begin{tabular}{lcc}
        \toprule
        Algorithm & Time Complexity & Space Complexity \\
        \midrule
        Binary Search & $O(\log n)$ & $O(1)$ \\
        Merge Sort    & $O(n\log n)$ & $O(n)$ \\
        Breadth-First Search & $O(V+E)$ & $O(V)$ \\
        \bottomrule
    \end{tabular}
\end{table}

Table~\ref{tab:complexity} compares the time and space
complexities of several algorithms.

\subsection{Full-Width Table}

The \texttt{tabularx} environment automatically adjusts column
widths to fit a specified total width.

\begin{table}[htbp]
    \centering
    \caption{Graph traversal algorithms.}
    \label{tab:traversal}

    \begin{tabularx}{\textwidth}{lXc}
        \toprule
        Algorithm & Description & Complexity \\
        \midrule
        BFS &
        Explores vertices level by level, starting from a source
        vertex. &
        $O(V+E)$ \\

        DFS &
        Explores each branch as deeply as possible before
        backtracking. &
        $O(V+E)$ \\
        \bottomrule
    \end{tabularx}
\end{table}

% ============================================================


% ============================================================
\section{Exercises}
% ============================================================

The \texttt{exercise} environment provides numbered exercises.

\begin{exercise}

Given an array of $n$ integers, find the maximum element.

What is the time complexity of a straightforward solution?

\end{exercise}

The exercise environment also accepts an optional title.

\begin{exercise}[Binary Search]

Given a sorted array of integers and a target value $x$,
implement binary search.

What is the time complexity of binary search?

\end{exercise}

Exercises can contain mathematics, lists, source code,
and other standard LaTeX content.

\begin{exercise}[BFS Implementation]

Given an unweighted graph $G=(V,E)$ and a starting vertex $s$,
implement breadth-first search.

A possible implementation is:

\begin{code}{python}
from collections import deque

def bfs(graph, start):
    visited = {start}
    queue = deque([start])

    while queue:
        u = queue.popleft()

        for v in graph[u]:
            if v not in visited:
                visited.add(v)
                queue.append(v)

    return visited
\end{code}

What is the time complexity of this implementation in terms of
$|V|$ and $|E|$?

\end{exercise}

\newpage
% ============================================================
\section{Information Boxes}
% ============================================================

The document system provides several semantic environments for
highlighting teaching information.

\subsection{Important Information}

\begin{important}

BFS computes shortest-path distances from the starting vertex
in an unweighted graph.

\end{important}

\subsection{Notes}

\begin{note}

When using an adjacency-list representation, iterating over all
neighbors of every vertex takes $O(|V|+|E|)$ time.

\end{note}

\subsection{Warnings}

\begin{warning}

Do not use ordinary BFS for weighted graphs when edge weights
can differ.

\end{warning}

% ============================================================
\section{Objectives and Tasks}
% ============================================================

\subsection{Objectives}

\begin{objective}

By the end of this section, students should be able to:

\begin{itemize}
    \item explain the BFS algorithm;
    \item implement BFS;
    \item analyze its time complexity;
    \item identify situations where BFS is appropriate.
\end{itemize}

\end{objective}

\subsection{Tasks}

A task without a title can be written as:

\begin{task}

Implement BFS using an adjacency-list representation.

\end{task}

A task can also have an optional title:

\begin{task}[Complexity Analysis]

Determine the time and space complexity of your implementation.

\end{task}

% ============================================================
\section{Questions, Answers, and Solutions}
% ============================================================

\subsection{Questions}

The \texttt{question} environment supports an optional label,
which can be used for marks or another short piece of information.

\begin{question}[2 points]

What is the time complexity of binary search?

\end{question}

\begin{answer}

The time complexity of binary search is

\[
    O(\log n).
\]

\end{answer}

\begin{question}[3 points]

Why does BFS find shortest paths in an unweighted graph?

\end{question}

\begin{solution}

BFS explores vertices in increasing order of their distance from
the starting vertex. Therefore, the first time a vertex is
visited, it has been reached through a shortest path.

\end{solution}

% ============================================================
\section{Algorithms}
% ============================================================

The \texttt{csdept-algorithms} package provides pseudocode through
\texttt{algorithm2e}.

\begin{algorithm}[H]
\caption{Breadth-First Search}

\KwIn{Graph $G=(V,E)$ and starting vertex $s$}
\KwOut{Distance array $d$}

\ForEach{$v \in V$}{
    $d[v] \gets \infty$\;
}

$d[s] \gets 0$\;
$Q \gets [s]$\;

\While{$Q$ is not empty}{
    $u \gets$ first element of $Q$\;
    remove $u$ from $Q$\;

    \ForEach{$v$ adjacent to $u$}{
        \If{$d[v] = \infty$}{
            $d[v] \gets d[u] + 1$\;
            append $v$ to $Q$\;
        }
    }
}

\KwRet{$d$}

\end{algorithm}

The algorithm package also provides custom inputs and outputs:

\begin{algorithm}[H]
\caption{Maximum Element}

\KwIn{Array $A$ containing $n$ elements}
\KwOut{The maximum element of $A$}

$maximum \gets A[1]$\;

\For{$i \gets 2$ \KwTo $n$}{
    \If{$A[i] > maximum$}{
        $maximum \gets A[i]$\;
    }
}

\KwRet{$maximum$}

\end{algorithm}



% ============================================================
% Graph Theory
% ============================================================

\section{Graph Theory}

The \texttt{csdept-graphs} package provides automatic graph
drawing through TikZ. Vertices and edges are specified directly,
while the layout algorithm determines their positions.

\subsection{Undirected Graph}

Use \texttt{--} to represent an undirected edge.

\begin{figure}[htbp]
    \centering

    \begin{tikzpicture}
        \graph [
            spring layout,
            nodes={
                draw,
                circle,
                minimum size=8mm
            },
            node distance=15mm
        ] {
            A -- B;
            A -- C;
            B -- C;
            B -- D;
            C -- D;
        };
    \end{tikzpicture}

    \caption{An undirected graph.}
    \label{fig:undirected-graph}
\end{figure}

\subsection{Directed Graph}

Use \texttt{->} to represent a directed edge.

\begin{figure}[htbp]
    \centering

    \begin{tikzpicture}
        \graph [
            layered layout,
            nodes={
                draw,
                circle,
                minimum size=8mm
            },
            edges={
                {Stealth-}
            }
        ] {
            A -> B;
            A -> C;
            B -> D;
            C -> D;
            D -> A;
        };
    \end{tikzpicture}

    \caption{A directed graph.}
    \label{fig:directed-graph}
\end{figure}

\subsection{Weighted Graph}

Edge labels can be specified using the quote syntax.
The same approach works for directed and undirected graphs.

\begin{figure}[htbp]
    \centering

    \begin{tikzpicture}
        \graph [
            spring layout,
            nodes={
                draw,
                circle,
                minimum size=8mm
            },
            edges={
                nodes={
                    font=\small,
                    fill=white,
                    inner sep=1pt
                }
            }
        ] {
            A --["4"] B;
            A --["2"] C;
            B --["7"] C;
            B --["3"] D;
            C --["5"] D;
        };
    \end{tikzpicture}

    \caption{An undirected weighted graph.}
    \label{fig:weighted-graph}
\end{figure}

\subsection{Self-Loops}

A self-loop is an edge that connects a vertex to itself.
It can be specified by using the same vertex as both endpoints.

\begin{figure}[htbp]
    \centering

    \begin{tikzpicture}
        \graph [
            spring layout,
            nodes={
                draw,
                circle,
                minimum size=10mm
            }
        ] {
            A ->[loop above] A;
            A -> B;
            B -> C;
            C ->[loop below] C;
        };
    \end{tikzpicture}

    \caption{A directed graph containing self-loops.}
    \label{fig:self-loops}
\end{figure}

\subsection{Directed Weighted Graph with a Self-Loop}

The following example combines directed edges, edge weights,
and a self-loop.

\begin{figure}[htbp]
    \centering

    \begin{tikzpicture}
        \graph [
            spring layout,
            nodes={
                draw,
                circle,
                minimum size=10mm
            },
            edges={
                nodes={
                    font=\small,
                    fill=white,
                    inner sep=1pt
                }
            }
        ] {
            A ->["3"] B;
            B ->["6"] C;
            C ->["2"] A;
            A ->["5", loop above] A;
        };
    \end{tikzpicture}

    \caption{A directed weighted graph with a self-loop.}
    \label{fig:combined-graph}
\end{figure}

% ============================================================

% ============================================================
\section{Source Code}
% ============================================================

The system uses \texttt{minted} and Pygments for syntax
highlighting.

No LaTeX escaping is required inside a \texttt{code} environment.
The contents are ordinary source code.

\subsection{Python}

\begin{code}{python}
def binary_search(a, x):
    left = 0
    right = len(a) - 1

    while left <= right:
        mid = (left + right) // 2

        if a[mid] == x:
            return mid
        elif a[mid] < x:
            left = mid + 1
        else:
            right = mid - 1

    return -1
\end{code}

\subsection{C++}

\begin{code}{cpp}
#include <iostream>
#include <vector>

int main() {
    std::vector<int> values = {1, 2, 3, 4, 5};

    for (int x : values) {
        std::cout << x << '\n';
    }

    return 0;
}
\end{code}

\subsection{Java}

\begin{code}{java}
public class Main {
    public static void main(String[] args) {
        System.out.println("Hello, world!");
    }
}
\end{code}

% ============================================================
\section{Lists and General LaTeX}
% ============================================================

Unordered lists work normally:

\begin{itemize}
    \item Graphs
    \item Trees
    \item Searching
    \item Sorting
\end{itemize}

Ordered lists work normally:

\begin{enumerate}
    \item Read the input.
    \item Process the graph.
    \item Produce the result.
\end{enumerate}

Nested lists are supported as well:

\begin{enumerate}
    \item Graph traversal

    \begin{itemize}
        \item BFS
        \item DFS
    \end{itemize}

    \item Shortest paths

    \begin{itemize}
        \item Dijkstra
        \item Bellman--Ford
    \end{itemize}
\end{enumerate}

% ============================================================
\section{Hyperlinks and References}
% ============================================================

Hyperlinks can be written using the standard \texttt{hyperref}
commands.

For example:

\begin{verbatim}
\url{https://www.example.com}
\end{verbatim}

A live hyperlink can also be written directly:

\url{https://www.example.com}

The document class configures hyperlink colors automatically.

% ============================================================
\section{Combined Example}
% ============================================================

The following example demonstrates how the different components
can be used together in an ordinary teaching document.

\begin{objective}

Students should understand how BFS explores an unweighted graph
and how its complexity depends on the graph representation.

\end{objective}

\begin{important}

For an adjacency-list representation, BFS runs in

\[
    O(|V|+|E|).
\]

\end{important}

\begin{exercise}[Graph Traversal]

Consider the graph $G=(V,E)$ where

\[
    V = \{A,B,C,D,E\}
\]

and

\[
    E =
    \{
        (A,B),
        (A,C),
        (B,D),
        (C,E)
    \}.
\]

\begin{task}[Traversal Order]

Starting from vertex $A$, determine the order in which BFS
visits the vertices.

\end{task}

\begin{question}[2 points]

What is the distance from $A$ to $E$?

\end{question}

\begin{answer}

The shortest path is

\[
    A \rightarrow C \rightarrow E,
\]

so the distance is $2$.

\end{answer}

\begin{code}{python}
from collections import deque

def bfs(graph, start):
    distance = {start: 0}
    queue = deque([start])

    while queue:
        u = queue.popleft()

        for v in graph[u]:
            if v not in distance:
                distance[v] = distance[u] + 1
                queue.append(v)

    return distance
\end{code}

\end{exercise}

% ============================================================
\section{Conclusion}
% ============================================================

This document demonstrates the main functionality provided by
the Computer Science Department document system.

A professor can create a document by loading the class:

\begin{verbatim}
\documentclass{csdept}
\end{verbatim}

and then defining the required metadata:

\begin{verbatim}
\course{Algorithms and Data Structures}
\teacher{Dr. Firstname Lastname}
\academicyear{2026--2027}
\documenttype{Tutorial Sheet}
\end{verbatim}

The professor can then freely combine ordinary LaTeX with the
department's mathematical, theorem, algorithm, source-code,
exercise, and teaching environments.

\end{document}

